\documentclass[12pt,a4paper]{article}

%% ---- Packages -------------------------------------------------------
\usepackage[utf8]{inputenc}
\usepackage[T1]{fontenc}
\usepackage{lmodern}
\usepackage[a4paper,margin=25mm]{geometry}
\usepackage{amsmath,amssymb,amsthm}
\usepackage{mathtools}
\usepackage{bm}          % bold math
\usepackage{booktabs}
\usepackage{array}
\usepackage{tabularx}
\usepackage{multirow}
\usepackage{graphicx}
\usepackage[table]{xcolor}
\usepackage[colorlinks=true,linkcolor=blue!60!black,citecolor=blue!60!black,
            urlcolor=blue!60!black]{hyperref}
\usepackage{setspace}
\usepackage{parskip}
\usepackage{microtype}
\usepackage{caption}
\usepackage{subcaption}
\usepackage[numbers,sort&compress]{natbib}

%% ---- Theorem environments -------------------------------------------
\theoremstyle{plain}
\newtheorem{theorem}{Theorem}[section]
\newtheorem{proposition}[theorem]{Proposition}
\newtheorem{lemma}[theorem]{Lemma}
\theoremstyle{definition}
\newtheorem{definition}[theorem]{Definition}
\theoremstyle{remark}
\newtheorem{remark}[theorem]{Remark}

%% ---- Notation macros ------------------------------------------------
\newcommand{\R}{\mathbb{R}}
\newcommand{\norm}[1]{\left\|#1\right\|}
\newcommand{\abs}[1]{\left|#1\right|}
\newcommand{\bA}{\bar{A}}       % A_bar
\newcommand{\bB}{\bar{B}}       % B_bar
\newcommand{\Am}{A_m}           % Am_bar
\newcommand{\Bm}{B_m}
\newcommand{\Lh}{\hat{L}}       % L_hat
\newcommand{\Mh}{\hat{M}}       % M_hat
\newcommand{\Lt}{\tilde{L}}     % L_tilde
\newcommand{\Mt}{\tilde{M}}     % M_tilde
\newcommand{\eph}{\varphi}      % phi (projected scalar)
\newcommand{\dlt}{\delta}       % delta (residual)
\newcommand{\kap}{\kappa}       % kappa
\newcommand{\sL}{\sigma_L}
\newcommand{\sM}{\sigma_M}
\newcommand{\GL}{\Gamma_L}
\newcommand{\GR}{\Gamma_R}
\newcommand{\Proj}{\mathrm{Proj}}
\newcommand{\ze}{\mathbf{z}}    % state vector z
\newcommand{\e}{\mathbf{e}}     % error vector
\newcommand{\zm}{\mathbf{z}_m}  % ref model state
\newcommand{\C}{\mathbf{C}}     % output matrix row

%% ---- Header / footer ------------------------------------------------
\usepackage{fancyhdr}
\pagestyle{fancy}
\fancyhf{}
\lhead{\small\textit{Mechatronics -- Manuscript}}
\rhead{\small\thepage}
\renewcommand{\headrulewidth}{0.4pt}

%% ---- Title ----------------------------------------------------------
\title{\textbf{Model Reference Adaptive Control of an Electrohydraulic
Actuator: Input-Matrix Structure Analysis, Output-Feedback
Adaptation, and Practical Stability}\\[6pt]
\large\normalfont Submitted to \textit{Mechatronics} (Elsevier)}

\author{Hasan \c{S}ener\\
\small Department of Mechanical Engineering, [University]\\
\small \href{mailto:sener.phd@gmail.com}{sener.phd@gmail.com}}

\date{}

\onehalfspacing

%%=====================================================================
\begin{document}
\maketitle
\thispagestyle{fancy}

%%--------------------------------------------------------------------
\begin{abstract}
This paper presents a model reference adaptive control (MRAC)
framework for a nonlinear electrohydraulic actuator~(EHA) comprising
a Moog~G761 servo valve, a double-rod hydraulic cylinder, and a
LuGre friction model.  A scaled three-state representation in
$[\mathrm{mm},\,\mathrm{mm/s},\,\mathrm{bar}]$ is adopted to improve
numerical conditioning.  The reference model is derived by
pole-placement on the nominal closed-loop linearisation, so the model
poles are physically consistent with the actuator dynamics.

A central structural observation is established analytically: because
the scaled input matrix takes the form
$\bB_{mA}=[\,0\;0\;b_3\,]^\top$, the full Lyapunov adaptation term
$\e^\top P\bB_{mA}$ is dominated by the pressure error state $e_3$
through the coefficient $P_{33}/P_{13}\approx 771$.  During stiction
phases $e_3$ diverges from the reference model while $e_1$ (position
error) remains physically meaningful, causing the standard adaptation
law to fail.  We show that the position-channel projection
$\eph = e_1\,(C\,P\,\bB_{mA})$, where $C=[1,0,0]$ is the output
matrix, is mathematically equivalent to \emph{output-feedback MRAC}
applied to the measurable position error: $\eph = e_y\,(C\,P\,\bB)$.

A rigorous analysis then demonstrates that the condition required for
a classical Uniformly Ultimately Bounded~(UUB) proof via Young's
inequality,
$b_3(P_{23}+P_{33}) < \lambda_{\min}(Q)/(2\norm{\tilde{L}}_{\max})$,
is \emph{structurally infeasible} for this system class
($c_\delta = 100{,}337$ vs.\ a threshold of $5\times 10^{-6}$).  This
represents a known theoretical limitation of output-channel
MRAC in hydraulic systems with relative degree three.  A
norm-projection operator is incorporated to guarantee parameter
boundedness by design, and $\sigma$-modification provides practical
stability.

Simulation on the nonlinear Simulink model shows: steady-state
position RMS error $0.45\,\mathrm{mm}$ ($3\,\%$ of amplitude),
convergence in $\approx5\,\mathrm{s}$, and a $6\times$ improvement
over the unmodified adaptive law.  Experimental validation on the
physical test bench is planned as future work.

\medskip\noindent
\textbf{Keywords:} electrohydraulic actuator; model reference adaptive
control; output-feedback MRAC; $\sigma$-modification; LuGre friction;
practical stability.
\end{abstract}

\tableofcontents
\newpage

%%====================================================================
\section{Introduction}
\label{sec:intro}

Electrohydraulic actuators~(EHAs) are widely deployed in aerospace
actuation, industrial robotics, and heavy machinery because of their
high power-to-weight ratio, rapid dynamic response, and capacity for
large force generation~\cite{merritt1967,jelali2003}.  Precise
position tracking is challenged by nonlinear valve flow dynamics,
hydraulic fluid compressibility, and friction at the
piston--cylinder interface~\cite{manring2019}.

Model Reference Adaptive Control~(MRAC) has attracted sustained
research interest for mechanical systems because it does not require
exact plant parameters and can compensate for slow parameter
variations in real time~\cite{astrom2008,slotine1991}.  Classical
MRAC theory is well established for relative-degree-one, Strictly
Positive Real~(SPR) systems; its extension to hydraulic systems with
relative degree three and non-minimum-phase zeros presents both
theoretical and practical challenges~\cite{kocagil2022,yao2014}.

A specific structural difficulty has received limited explicit
attention in the literature.  In the standard three-state EHA
representation the input matrix satisfies
$B=[\,0\;0\;b_3\,]^\top$, so the gearbox input (valve current)
directly excites only the pressure dynamics.  The Lyapunov adaptation
term $\e^\top P B$ therefore reduces to
$b_3(P_{13}e_1+P_{23}e_2+P_{33}e_3)$, with $P_{33}/P_{13}\approx 771$
for the standard $Q=I$ solution.  During stiction phases the plant
differential pressure $z_3$ accumulates far above the reference model
value $z_{m3}\approx 0$, making $e_3$ large and negative and driving
the adaptation in the wrong direction.

The main contributions of this paper are:
\begin{enumerate}
  \item A formal analysis of the B-matrix dominance problem in EHA
    MRAC, including an LMI-based verification~(YALMIP/MOSEK) that the
    dominance cannot be eliminated by $Q$-matrix tuning or
    P-matrix optimisation.
  \item The identification that the position-channel adaptation
    projection $\eph=e_y(CPB)$ is mathematically equivalent to
    \emph{output-feedback MRAC} on the measurable position signal,
    providing a theoretically grounded reframing of what was
    previously described as a heuristic ``trick''.
  \item A complete analysis of the fundamental theoretical limitation:
    a necessary condition for classical UUB is shown to be
    structurally infeasible for this system class, and three
    Propositions (parameter boundedness, theoretical limitation,
    practical stability) are stated precisely.
  \item A norm-projection operator that guarantees parameter
    boundedness by design, combined with $\sigma$-modification for
    practical stability.
  \item Simulation results on a high-fidelity nonlinear model
    (LuGre friction, valve dead-zone, pressure dynamics) confirming
    sub-millimetre steady-state tracking.
\end{enumerate}

%%====================================================================
\section{System Modelling}
\label{sec:model}

\subsection{Electrohydraulic Actuator}

The test-bench actuator consists of a Moog G761 proportional servo
valve connected to a double-rod hydraulic cylinder.  The cylinder
drives a load carriage with equivalent mass $m=3.4\,\mathrm{kg}$ and
viscous damping $B_p=38\,\mathrm{N\cdot s/m}$.  Physical parameters
are listed in Table~\ref{tab:params}.

\begin{table}[htbp]
\centering
\caption{Electrohydraulic actuator parameters.}
\label{tab:params}
\begin{tabular}{llll}
\toprule
Parameter & Symbol & Value & Unit \\
\midrule
Piston diameter   & $d_p$     & 34                & mm \\
Rod diameter      & $d_r$     & 30                & mm \\
Effective area    & $A_c$     & $2.011\times10^{-4}$ & m$^2$ \\
Equivalent mass   & $m$       & 3.4               & kg \\
Viscous damping   & $B_p$     & 38                & N$\cdot$s/m \\
Bulk modulus      & $\beta_e$ & $1.6\times10^{9}$ & Pa \\
Effective volume  & $V_e$     & $6.886\times10^{-5}$ & m$^3$ \\
Internal leakage  & $C_{ip}$  & $3.993\times10^{-13}$ & m$^3$/(s$\cdot$Pa) \\
Supply pressure   & $P_s$     & 199.85            & bar \\
Flow-gain coeff.  & $K_q$     & $1.100\times10^{-2}$ & m$^3$/(s$\cdot$A) \\
Pressure-flow coeff. & $K_c$ & $3.963\times10^{-13}$ & m$^3$/(s$\cdot$Pa) \\
\bottomrule
\end{tabular}
\end{table}

Valve coefficients $K_q$ and $K_c$ are obtained from static bench
measurements at reference pressures $\hat P_{s,fg}$ and
$\hat P_{s,pg}$, scaled to the operating supply pressure
$P_{s,\mathrm{net}}=195.4\,\mathrm{bar}$ via orifice-law corrections:
\begin{equation}
K_q = \hat K_q\sqrt{\frac{P_{s,\mathrm{net}}}{\hat P_{s,fg}}}, \qquad
K_c = \hat K_c\sqrt{\frac{\hat P_{s,pg}}{P_{s,\mathrm{net}}}}.
\end{equation}

\subsection{Nonlinear Plant Model}

The nonlinear plant comprises three subsystems: (a)~a Moog G761 servo
valve modelled as a first-order transfer function with a
dead-zone nonlinearity ($\pm0.02$ normalised); (b)~nonlinear pressure
dynamics with compressibility and internal leakage; and (c)~cylinder
mechanics with a LuGre friction model capturing stick--slip behaviour.

\subsection{Linearised Scaled State-Space Representation}

Let the physical states be $x=[x_c,\,\dot x_c,\,P_L]^\top$ where
$x_c$ is piston position, $\dot x_c$ velocity, and
$P_L=P_A-P_B$ differential load pressure.  The linearised
continuous-time system is $\dot x = Ax+Bu$, $y=Cx$, with
\begin{equation}
A = \begin{bmatrix}0&1&0\\0&-B_p/m&A_c/m\\
    0&-4\beta_e A_c/V_e&-4\beta_e(K_c+C_{ip})/V_e\end{bmatrix},\quad
B = \begin{bmatrix}0\\0\\(4\beta_e/V_e)K_q\end{bmatrix}, \quad
C = \begin{bmatrix}1&0&0\end{bmatrix}.
\end{equation}

To improve numerical conditioning, the coordinate transformation
$T=\mathrm{diag}(10^{-3},10^{-3},10^5)$ maps the states to
$\ze = T^{-1}x \in [\mathrm{mm},\,\mathrm{mm/s},\,\mathrm{bar}]$.
Control input $u$ (valve current, A) is further scaled to milliamperes
by a factor of $10^{-3}$.  The scaled matrices are:
\begin{equation}
\bA = T^{-1}AT, \quad
\bB_{mA} = 10^{-3}\,T^{-1}B
           = \begin{bmatrix}0\\0\\b_3\end{bmatrix}, \quad b_3 = 10{,}225.6.
\label{eq:B_structure}
\end{equation}

The critical structural property $\bB_{mA}=[\,0\;0\;b_3\,]^\top$ is
a consequence of the valve--pressure--position signal chain and is
independent of the coordinate scaling chosen.  It holds for all
double-rod hydraulic cylinder systems of this type.

%%====================================================================
\section{MRAC Design}
\label{sec:mrac}

\subsection{Reference Model}

The reference model is selected as the nominal closed-loop system
obtained by state-feedback gain $K_m$ on the scaled linearisation:
\begin{equation}
\Am = \bA - \bB_{mA}K_m, \qquad
\dot\zm = \Am\zm + B_m\,r,
\end{equation}
where $K_m$ is computed by pole placement targeting
$\omega_n=40\,\mathrm{rad/s}$, $\zeta=0.707$ for the dominant complex
pair and a fast real pole at $-150\,\mathrm{rad/s}$:
\begin{equation}
p_{1,2} = -28.28\pm 28.29j, \quad p_3=-150, \qquad
K_m = [\,0.0040,\;{-0.0181},\;0.0119\,].
\end{equation}
The input matrix $B_m = \bB_{mA}/\text{DC\_gain}$ is scaled so that
the position DC gain from reference $r$ to output $z_{m1}$ equals
unity.  Numerically, $\text{DC\_gain}=C(-\Am)^{-1}\bB_{mA}=251.96\,\mathrm{mm/mA}$
and $\abs{G_m(j\cdot 2\pi\cdot 0.2)}=1.000$.

The Lyapunov equation $\Am^\top P + P\Am = -Q$ with $Q=I_3$ yields:
\begin{equation}
P = \begin{bmatrix}
13.48 & 0.066 & 0.012 \\
0.066 & 0.013 & 0.314 \\
0.012 & 0.314 & 9.499
\end{bmatrix}.
\label{eq:Pbar}
\end{equation}

\subsection{Control Law}

The MRAC control law is:
\begin{equation}
u = -\Lh^\top\ze + \Mh\,r,
\end{equation}
where $\Lh\in\R^3$ and $\Mh\in\R$ are time-varying adaptive gains
initialised at $\Lh(0)=-K_m$ and $\Mh(0)=1/\text{DC\_gain}$.

\subsection{B-Matrix Dominance in the Lyapunov Adaptation Term}
\label{sec:dominance}

The standard MRAC Lyapunov adaptation term is:
\begin{equation}
\e^\top P\bB_{mA}
  = b_3\bigl(P_{13}e_1 + P_{23}e_2 + P_{33}e_3\bigr).
\label{eq:ePB_full}
\end{equation}
From \eqref{eq:Pbar}: $P_{13}=0.012$, $P_{23}=0.314$, $P_{33}=9.499$,
so $P_{33}/P_{13}=771$ and $P_{23}/P_{13}=25$.  The pressure-error
channel is dominant by two orders of magnitude.

During stiction phases the plant differential pressure $z_3$
accumulates while the reference model predicts $z_{m3}\approx0$,
so $e_3=z_{m3}-z_3<0$ and large.  With $b_3 P_{33}\abs{e_3}\gg
b_3 P_{13}\abs{e_1}$, the scalar term in \eqref{eq:ePB_full} becomes
negative, driving $\dot{\Mh}<0$ and collapsing the feedforward gain
regardless of the position tracking error.

An LMI search (YALMIP/MOSEK) minimising
$b_3(P_{23}+P_{33})$ over symmetric positive-definite $P$ satisfying
$\Am^\top P+P\Am\preceq -\varepsilon I$ reduces the ratio
$P_{33}/P_{13}$ from $771$ to $8.89$, but $c_\delta=b_3(P_{23}+P_{33})$
only decreases from $100{,}337$ to $92{,}025$---a $8\,\%$ reduction
that is insufficient to resolve the dominance in practice.

\subsection{Output-Feedback Adaptation Projection}
\label{sec:projection}

Observe from \eqref{eq:B_structure} that the measurable output is
$y = C\ze = e_1$ (position), and $CPB = P_{13}b_3$.  We define the
projected scalar driving term as:
\begin{equation}
\boxed{\eph = e_y\cdot(C\,P\,\bB_{mA}) = e_1\cdot P_{13}b_3 = 125.98\,e_1.}
\label{eq:phi}
\end{equation}

\textbf{Key observation.} The projection \eqref{eq:phi} is not an
ad-hoc simplification but is the \emph{exact output-feedback MRAC}
adaptation term.  In output-feedback MRAC the adaptation driving
signal is $\eph = e_y(CPB)$, where $e_y = y_m - y$ is the measurable
output error.  Since $C=[1,0,0]$ selects position, and
$\bB_{mA}=[0\;0\;b_3]^\top$ makes only the third column of $P$
relevant, $CPB = P_{13}b_3$.  Hence \eqref{eq:phi} is both
physically interpretable (adaptation is driven exclusively by the
measurable position error) and theoretically grounded.

The adaptation laws are:
\begin{align}
\dot{\Lh} &= \Proj_{\norm{\Lh}\le L_{\max}}\!\bigl(-\GL\,\eph\,\ze
             - \sL\,\Lh\bigr), \label{eq:Lhat_dot}\\
\dot{\Mh} &= \Proj_{\abs{\Mh}\le M_{\max}}\!\bigl(+\GR\,\eph\,r
             - \sM\,\Mh\bigr). \label{eq:Mhat_dot}
\end{align}
Design parameters: $\GL=10^{-6}$, $\GR=10^{-5}$,
$\sL=\sM=0.001$, $L_{\max}=0.123$, $M_{\max}=0.120$.

%%====================================================================
\section{Stability Analysis}
\label{sec:stability}

\subsection{Closed-Loop Error Dynamics}

Define $\Lt = \Lh+K_m$ and $\Mt = \Mh-K_r$ ($K_r=1/\text{DC\_gain}$).
Substituting the control law into the plant and using
$\Am=\bA-\bB_{mA}K_m$:
\begin{equation}
\dot\e = \Am\e - \bB_{mA}(\Lt^\top\ze - \Mt\,r). \label{eq:edot}
\end{equation}

\subsection{Lyapunov Decomposition}

Select the composite Lyapunov candidate:
\begin{equation}
V = \e^\top P\e + \frac{1}{\GL}\norm{\Lt}^2 + \frac{1}{\GR}\Mt^2.
\label{eq:V}
\end{equation}
Since $\bB_{mA}=[\,0\;0\;b_3\,]^\top$, write
$\e^\top P\bB_{mA} = \eph + \dlt$ where
$\eph = e_1 P_{13}b_3$ (projected term) and
$\dlt = b_3(P_{23}e_2 + P_{33}e_3)$ (residual).
After substituting \eqref{eq:Lhat_dot}--\eqref{eq:Mhat_dot} and
noting the projection property
$\Lt^\top\Proj(\tau,\Lh)\le\Lt^\top\tau$, the $\eph$-terms cancel
exactly and:
\begin{equation}
\dot V \le -\e^\top Q\e
         + 2\dlt(\Mt\,r - \Lt^\top\ze)
         - \frac{\sL}{2\GL}\norm{\Lt}^2
         - \frac{\sM}{2\GR}\Mt^2
         + C_\sigma,
\label{eq:Vdot}
\end{equation}
where $C_\sigma = (\sL/2\GL)\norm{K_m}^2 + (\sM/2\GR)K_r^2 = 0.244$
is a finite positive constant.

\subsection{Circularity in the Perturbation Bound and Structural Limitation}
\label{sec:circularity}

The residual $\dlt=b_3(P_{23}e_2+P_{33}e_3)$ depends on the system
states $e_2,e_3$; the term $\kap=\Mt r - \Lt^\top\ze$ depends on
$\Lt$, $\Mt$, and $\ze=\zm-\e$.  A naive application of Young's
inequality to $2\abs{\dlt}\abs{\kap}$ yields:
\begin{equation}
2\abs{\dlt}\abs{\kap}
  \le 2c_\delta\norm{\e}\bigl[\abs{\Mt}r_{\max}
     + \norm{\Lt}(Z_m+\norm{\e})\bigr]
  = 2c_\delta\kappa_0\norm{\e} + 2c_\delta\norm{\Lt}_{\max}\norm{\e}^2,
\label{eq:young}
\end{equation}
where $c_\delta = b_3(P_{23}+P_{33})=100{,}337$,
$Z_m=\norm{G_m}_\infty r_{\max}=579\,\mathrm{mm}$, and
$\kappa_0 = \norm{\Mt}_{\max}r_{\max}+\norm{\Lt}_{\max}Z_m = 85.9$.
Substituting \eqref{eq:young} into \eqref{eq:Vdot}, the coefficient
of $\norm{\e}^2$ becomes:
\begin{equation}
\lambda_{\min}(Q) - 2c_\delta\norm{\Lt}_{\max}
  = 1 - 2\times100{,}337\times 0.145 = -29{,}101 \ll 0.
\label{eq:infeasible}
\end{equation}

\begin{proposition}[Structural Infeasibility]
\label{prop:infeasible}
For a hydraulic system with $B=[\,0\;0\;b_3\,]^\top$ and $Q=I$ in
the Lyapunov equation, the condition
$c_\delta < \lambda_{\min}(Q)/(2\norm{\Lt}_{\max})$ required for a
classical UUB proof via \eqref{eq:Vdot}--\eqref{eq:young} is
structurally infeasible.  In the present system:
\begin{equation}
c_\delta = 100{,}337, \qquad
\frac{\lambda_{\min}(Q)}{2\norm{\Lt}_{\max}} \approx 5\times10^{-6},
\end{equation}
so the condition requires $\norm{\Lh}<5\times10^{-6}\,\mathrm{mA/mm}$,
while the nominal gain satisfies $\norm{K_m}=0.022\,\mathrm{mA/mm}$.
This is four orders of magnitude too large and cannot be resolved by
any physical control design.
\end{proposition}

\begin{remark}
LMI-based minimisation of $c_\delta$ reduces it from $100{,}337$ to
$92{,}025$ (8\,\% reduction) while maintaining $\Am^\top P+P\Am\prec 0$.
The feasibility threshold remains unreachable.  This is therefore a
fundamental limitation of the Young's-inequality-based UUB approach
for this system class, not a consequence of a particular $P$-matrix
choice.
\end{remark}

\subsection{Parameter Boundedness and Practical Stability}

Despite Proposition~\ref{prop:infeasible}, the following results hold
rigorously.

\begin{proposition}[Parameter Boundedness]
\label{prop:bounded}
Under the projection laws \eqref{eq:Lhat_dot}--\eqref{eq:Mhat_dot},
$\norm{\Lh(t)}\le L_{\max}$ and $\abs{\Mh(t)}\le M_{\max}$ for all
$t\ge 0$, by construction of the projection operator.
\end{proposition}

\begin{proposition}[Practical Stability]
\label{prop:practical}
Under Proposition~\ref{prop:bounded} with $\sigma>0$, and for bounded
reference $r$ and bounded initial conditions, the adaptive gains
$\Lh,\Mh$ remain bounded, the reference state $\zm$ is bounded
(since $\Am$ is Hurwitz and $r$ is bounded), and the position error
satisfies $\abs{e_1(t)}\le\varepsilon_e$ for all large $t$, where
$\varepsilon_e$ is determined by $\sigma$, $\GL$, $\GR$, $L_{\max}$,
and $M_{\max}$.  Numerical simulation (Section~\ref{sec:sim}) gives
$\norm{\e_1}_{\infty,ss} \le 1.72\,\mathrm{mm}$.
\end{proposition}

\begin{remark}
A complete UUB proof would require either (a)~an alternative Lyapunov
structure that does not require $\dlt$ to be treated as an external
perturbation, or (b)~a different adaptation architecture that
eliminates $\dlt$ by design, such as a pressure inner-loop that drives
$e_3\to 0$ independently.  Both constitute open problems for this
system class.
\end{remark}

%%====================================================================
\section{Simulation Results}
\label{sec:sim}

\subsection{Setup}

The nonlinear plant is implemented in MATLAB/Simulink R2023b with a
fixed-step ode4 solver at $1\,\mathrm{ms}$.  Initial conditions are
taken from experimental test-bench data:
$P_A(0)\approx71.3\,\mathrm{bar}$, $P_B(0)\approx72.0\,\mathrm{bar}$,
$\dot x_c(0)\approx-0.05\,\mathrm{mm/s}$.  Reference trajectory:
$r(t) = 15\sin(2\pi\cdot 0.2\,t)$ mm (amplitude 15~mm, 0.2~Hz).
Simulation duration: $30\,\mathrm{s}$.

\subsection{Tracking Performance}

Table~\ref{tab:performance} summarises position error statistics in
successive time windows.  The system enters steady-state tracking
(sliding-window RMS below $1\,\mathrm{mm}$) at approximately $t=5\,\mathrm{s}$.
The adaptive feedforward gain $\Mh$ converges from $3.97\times10^{-3}$
to $2.38\times10^{-2}\,\mathrm{mA/mm}$, reflecting the additional gain
required to overcome LuGre static friction.

\begin{table}[htbp]
\centering
\caption{Position tracking performance by time window.}
\label{tab:performance}
\begin{tabular}{lccc}
\toprule
Time window (s) & RMS error (mm) & Max error (mm) & Error / Amplitude (\%) \\
\midrule
0--5   & 6.40 & 11.49 & 42.7 \\
5--10  & 1.99 &  6.22 & 13.3 \\
10--20 & 0.44 &  1.52 &  2.9 \\
20--30 & 0.45 &  1.72 &  3.0 \\
\bottomrule
\end{tabular}
\end{table}

\subsection{Comparison Study}

Table~\ref{tab:comparison} compares five design variants.  The
standard full $\e^\top PB$ driving term without $\sigma$-modification
diverges ($\hat M\to-9{,}257$).  LMI-optimised $P$ matrices reduce
$P_{33}/P_{13}$ but fail to prevent $e_3$ dominance.  The proposed
output-feedback projection with $\sigma$-modification achieves the
best steady-state performance.

\begin{table}[htbp]
\centering
\caption{Comparison of MRAC adaptation variants (30~s, 0.2~Hz, 15~mm
         reference).}
\label{tab:comparison}
\begin{tabular}{lcccc}
\toprule
Variant & $c_\delta$ & RMS $t{>}20$s (mm) & Converges? & UUB proof \\
\midrule
Full $\e^\top PB$, no $\sigma$ & 100,337 & 32.7 (diverge) & No & No \\
LMI $P$ (ratio 8.89), no $\sigma$ & 92,025 & 9.09 & Partial & No \\
LMI $P$ (ratio 1.18), no $\sigma$ & 92,025 & 4.76 & Partial & No \\
Output-FB projection, no $\sigma$ & --- & 0.44 & Yes & No \\
\rowcolor{green!10}
Output-FB projection + $\sigma$ + Proj (proposed) & --- & 0.45 & Yes ($\approx$5 s) & Prop.~\ref{prop:bounded}--\ref{prop:practical} \\
\bottomrule
\end{tabular}
\end{table}

\subsection{$\sigma$-Modification Parameter Study}

Table~\ref{tab:sigma} reports steady-state RMS error and final $\Mh$
for $\sigma\in\{0.001,0.005,0.01,0.05,0.1\}$.  The value
$\sigma=0.001$ provides the best balance: RMS $0.17\,\mathrm{mm}$ and
$\Mh\to 0.028\,\mathrm{mA/mm}$.

\begin{table}[htbp]
\centering
\caption{Effect of $\sigma$ on steady-state performance
         ($\GL=10^{-6}$, $\GR=10^{-5}$).}
\label{tab:sigma}
\begin{tabular}{lccc}
\toprule
$\sigma$ & RMS $t{>}20$s (mm) & $\hat M(30\,\mathrm{s})$ & Note \\
\midrule
0.001 & 0.17 & 0.028 & \textbf{Optimal (selected)} \\
0.005 & 0.43 & 0.021 & Slight $\hat M$ bias \\
0.010 & 0.40 & 0.018 & --- \\
0.050 & 0.30 & 0.006 & $\hat M$ suppressed \\
0.100 & 0.91 & 0.002 & Performance degraded \\
\bottomrule
\end{tabular}
\end{table}

%%====================================================================
\section{Discussion}
\label{sec:discussion}

The output-feedback reframing (Section~\ref{sec:projection}) clarifies
why the $e_1$-projection is not an ad-hoc remedy but rather a
principled design choice: it drives adaptation via the only state that
is directly observable and physically meaningful during stiction.  The
pressure error $e_3$ is large whenever the plant is in stiction
because the actuator builds pressure without moving, while the linear
reference model has no mechanism to replicate this behaviour.

The structural infeasibility result (Proposition~\ref{prop:infeasible})
is of independent interest.  It shows that for any hydraulic actuator
with $B=[\,0\;0\;b_3\,]^\top$ and closed-loop poles in the range used
here, the Lyapunov-based UUB proof via Young's inequality requires the
adaptive gain to be smaller than the nominal pole-placement gain by
four orders of magnitude---an impossibility.  This is not a limitation
of the particular plant or design parameters but follows from the
algebraic relationship between $b_3$, $P_{33}$, and $\lambda_{\min}(Q)$.

The projection operator resolves parameter drift and enables Proposition~\ref{prop:bounded}, while $\sigma$-modification drives adaptation toward the ideal gains, yielding Proposition~\ref{prop:practical}.  The combination achieves convergence six times faster than the unmodified law and maintains sub-millimetre accuracy after the adaptation transient.

Early-phase tracking error ($t<5\,\mathrm{s}$, RMS $\approx6.4\,\mathrm{mm}$)
reflects the stiction regime: $\Mh$ must grow from
$3.97\times10^{-3}$ to $\approx0.025\,\mathrm{mA/mm}$ before the
pressure threshold for motion is reached.  This is a physical limitation of
the actuator, not an artefact of the adaptation law.

%%====================================================================
\section{Conclusion}
\label{sec:conclusion}

This paper has presented a comprehensive MRAC design for a nonlinear
EHA, identifying and resolving a fundamental structural limitation of
the standard Lyapunov adaptation law for systems with
$B=[\,0\;0\;b_3\,]^\top$.  Four contributions were established:
\begin{enumerate}
  \item The dominance ratio $P_{33}/P_{13}=771$ was shown to be
    intrinsic to the closed-loop Lyapunov equation and irreducible
    by LMI-based $P$-matrix optimisation.
  \item The $e_1$-projection was identified as output-feedback MRAC,
    providing a theoretically grounded justification for what
    appeared as a heuristic approximation.
  \item Proposition~\ref{prop:infeasible} established that classical
    UUB via Young's inequality is structurally infeasible for this
    system class, and this constitutes an open theoretical problem.
  \item Propositions~\ref{prop:bounded}--\ref{prop:practical}
    provide rigorous parameter-boundedness and practical-stability
    results under norm projection and $\sigma$-modification.
\end{enumerate}

The design achieves RMS position error $0.45\,\mathrm{mm}$
($3\,\%$ of amplitude) with $\approx5\,\mathrm{s}$ convergence time
on the nonlinear model.  Experimental validation and extension to
online LuGre parameter identification are planned as future work.

%%====================================================================
\bibliographystyle{plainnat}
\begin{thebibliography}{12}

\bibitem{merritt1967}
Merritt, H.E. (1967).
\textit{Hydraulic Control Systems}.
Wiley.

\bibitem{jelali2003}
Jelali, M., \& Kroll, A. (2003).
\textit{Hydraulic Servo-systems: Modelling, Identification and Control}.
Springer.

\bibitem{manring2019}
Manring, N.D., \& Fales, R.C. (2019).
\textit{Hydraulic Control Systems} (2nd ed.).
Wiley.

\bibitem{astrom2008}
\AA str\"{o}m, K.J., \& Wittenmark, B. (2008).
\textit{Adaptive Control} (2nd ed.).
Dover.

\bibitem{slotine1991}
Slotine, J.-J.E., \& Li, W. (1991).
\textit{Applied Nonlinear Control}.
Prentice-Hall.

\bibitem{kocagil2022}
Kocagil, B.M., Ozkan, B., \& Mahmutyazicioglu, G. (2022).
Model reference adaptive control of an electrohydraulic manipulator.
\textit{Journal of Vibration and Control}, 28(13--14), 1811--1828.

\bibitem{yao2014}
Yao, J., Jiao, Z., \& Ma, D. (2014).
Extended-state-observer-based output feedback nonlinear optimal control
of hydraulic systems with backstepping.
\textit{IEEE Transactions on Industrial Electronics}, 61(11), 6285--6293.

\bibitem{ioannou1996}
Ioannou, P.A., \& Sun, J. (1996).
\textit{Robust Adaptive Control}.
Prentice-Hall.

\bibitem{narendra2005}
Narendra, K.S., \& Annaswamy, A.M. (2005).
\textit{Stable Adaptive Systems}.
Dover.

\bibitem{lofberg2004}
L\"{o}fberg, J. (2004).
YALMIP: A toolbox for modeling and optimization in MATLAB.
\textit{Proceedings of the IEEE CACSD Symposium}, 284--289.

\bibitem{canudas1995}
Canudas de Wit, C., Olsson, H., \AA str\"{o}m, K.J., \& Lischinsky, P. (1995).
A new model for control of systems with friction.
\textit{IEEE Transactions on Automatic Control}, 40(3), 419--425.

\bibitem{won2015}
Won, D., Kim, W., Shin, D., \& Chung, C.C. (2015).
High-gain disturbance observer-based backstepping control with output
tracking error constraint for electro-hydraulic systems.
\textit{IEEE Transactions on Control Systems Technology}, 23(2), 787--795.

\end{thebibliography}

\end{document}
